\section{Contribuciones tres y cuatro: Self-Play RL y Value-Based Filtering}

La sección anterior describe una sola ronda de planning; Cicero también necesita mejorar el modelo de políticas durante el entrenamiento y hacer una última revisión a nivel de comportamiento antes de enviar un mensaje. El self-play RL lleva a la policy más allá del nivel promedio de la imitation; el value-based filtering, por su parte, incorpora a la selección de candidatos la pregunta «¿cómo hará actuar esta frase a los demás?».

\subsection{Self-play RL: superar la imitación sin separarse de los humanos}

La slide 36 resalta la tercera contribución. El equipo no sustituye directamente el human model por un unconstrained self-play, sino que sigue usando piKL, de modo que el RL mejore en valor sin perder capacidad de explicar la distribución de las políticas humanas. La conclusión que se presentó en clase es que este enfoque también modela el comportamiento humano mejor que la imitation sola.

\lecturefigure{slide-36-main-contribution-selfplay.jpg}{Main contribution: mejorar la política con self-play RL regularizado con piKL}{Diapositiva V036 recuperada de la grabación oficial; Stanford CS25 Lecture 14}

\begin{knowledgebox}{Diferencia entre el RL en entrenamiento y la search en inferencia}
El self-play RL actualiza los parámetros para que la base policy mejore en muchas situaciones futuras; la inference-time search no necesariamente modifica los pesos, sino que invierte cómputo adicional de forma temporal en la situación actual. Cicero necesita ambas: el entrenamiento proporciona un prior utilizable y la inferencia lo refina según el diálogo y la situación concretos.
\end{knowledgebox}

\subsection{Por qué el filtrado del lenguaje debe considerar el strategic value}

La cuarta contribución parte de los candidatos ya generados. Los filtros habituales de toxicity, repetición y reglas solo detectan problemas superficiales; Cicero también se enfrenta a mensajes gramaticalmente impecables que arruinan su propia estrategia. Por ejemplo, el candidato absurdo que mostró el ponente amenaza con «sacarte del tablero» y, al mismo tiempo, le pide a la otra parte un «croissant». Las fallas más sutiles consisten en revelar sin querer un plan de ataque o en inducir a la otra parte a tomar acciones que perjudican al propio jugador.

\lecturefigure{slide-37-main-contribution-filtering.jpg}{Main contribution: filtrar los mensajes candidatos por estrategia y coherencia}{Diapositiva V037 recuperada de la grabación oficial; Stanford CS25 Lecture 14}

\subsection{La prueba causal del value-based filtering}

Sea $m$ el mensaje candidato, $a_{-i}$ las acciones futuras del receptor y $a_i$ las acciones planeadas de Cicero. El sistema puede usar el modelo de comportamiento para predecir la política de la otra parte antes y después del envío del mensaje, y comparar el valor del plan:
\begin{equation}
V(m)
=\mathbb{E}_{a_{-i}\sim p(\cdot\mid s,h,d,m)}
\left[Q_i(s,a_i,a_{-i})\right].
\end{equation}
Si $V(m)$ es claramente menor que el valor de no enviar nada o de enviar un mensaje de referencia, el candidato debe rechazarse. Lo que se evalúa aquí no es el sentimiento de la frase, sino el efecto causal que tiene sobre la ganancia futura al modificar el comportamiento de los demás.

\lecturefigure{slide-38-value-filtering.jpg}{Value-based filtering: predecir el efecto del mensaje sobre las acciones de los demás y la ganancia del plan}{Diapositiva V038 recuperada de la grabación oficial; Stanford CS25 Lecture 14}

\begin{knowledgebox}{Lectura de la figura: entradas y salidas del filtro}
A la izquierda está la generación de candidatos; en el centro, el action predictor simula la policy del receptor después de leer el candidato; a la derecha, el evaluator estima el valor del plan original de Cicero ante esa respuesta. Un candidato puede pasar las revisiones de gramática y seguridad y, aun así, descartarse en esta capa porque reduce el expected value. Por eso el filtrado es una counterfactual evaluation ligera y no una lista negra de palabras clave.
\end{knowledgebox}

\begin{warningbox}{Los errores del evaluator se convierten en un punto ciego sistemático}
Si el action predictor subestima el enojo, la represalia o la memoria de largo plazo de las personas, $V(m)$ quedará sobrestimado. El filtro solo es tan bueno como el modelo interno del mundo; no puede compensar de la nada las consecuencias sociales que el modelo no representa. Ese es justamente el origen de las limitaciones de truthfulness que se tratan más adelante.
\end{warningbox}

\begin{lstlisting}[caption={Cicero-style strategic dialogue closed loop},label={lst:cicero-loop}]
while game_is_active(state):
    human_anchor = predict_human_policies(state, history, dialogue)
    policy = pikl_search(state, human_anchor, value_model)
    own_plan, expected_others = choose_joint_intents(policy)

    candidates = dialogue_model.generate(
        state=state,
        history=history,
        dialogue=dialogue,
        own_intent=own_plan,
        recipient_intent=expected_others,
    )
    candidates = filter_nonsense_and_inconsistency(candidates)
    message = max(candidates, key=predicted_strategic_value)
    send(message)
    state, history, dialogue = observe_next_event()
\end{lstlisting}

Lo más importante de este pseudocódigo es la última línea: después de observar un mensaje nuevo o un estado nuevo, el sistema vuelve a predecir y a planear. Si el mensaje de salida se toma como el final del pipeline, se pierde la definición de agent; un verdadero agente debe hacer que los resultados de la comunicación regresen a la belief y a la policy.

\subsection{Resumen de la sección}

El self-play RL mejora la base policy a largo plazo, piKL conserva la capacidad de coordinarse con humanos y el value-based filtering revisa, antes de cada envío, las consecuencias estratégicas del lenguaje. Las tres muestran en conjunto que los candidatos en lenguaje natural deben pasar por un modelo de comportamiento y una función de valor, en lugar de ir directamente del decoder al entorno.
